\begin{solapartat}
Calculem la freqüència:
\[ f = \frac{E}{h}\ \text{\punts{0,2}} = 10\ \mathrm{eV}\,\frac{1,602\cdot 10^{-19}\ \mathrm{J}}{1\ \mathrm{eV}}\times\frac{1}{6,626\cdot 10^{-34}\ \mathrm{J\,s}}\ \text{\punts{0,1}} = 2,418\cdot 10^{15}\ \mathrm{Hz}\ \text{\punts{0,1}} \]

La longitud d'ona serà:
\[ \lambda = \frac{\mathrm{c}}{f}\ \text{\punts{0,2}} = 1,241\cdot 10^{-7}\ \mathrm{m}\times\frac{1\mathrm{nm}}{10^{-9}\mathrm{m}} = 124,1\ \mathrm{nm}\ \text{\punts{0,1}} \]

Com que el camp elèctric és constant tindrem:
\[ E = \frac{\Delta\mathrm{V}}{d}\ \text{\punts{0,2}} = 10\ \mathrm{N/C}\ \text{\punts{0,1}} \]
\end{solapartat}

\begin{solapartat}
Per trobar l'energia cinètica amb què surten els electrons des de l'ànode A, farem servir el principi de conservació de l'energia total:
\[ E_c^A + E_p^A = E_c^B + E_p^B\ \text{\punts{0,2}} \Rightarrow \]
\[ E_c^A = E_c^B + E_p^B - E_p^A = E_c^B + q_e(V_B - V_A)\ \text{\punts{0,2}} = 2,3\ \mathrm{eV} - 1\mathrm{e}\,(-3\mathrm{V})\ \text{\punts{0,2}} = 5,3\ \mathrm{eV}\ \text{\punts{0,2}} \]
Per trobar el treball d'extracció només caldrà que restem a l'energia dels fotons, la energia cinètica dels electrons emesos:
\[ W = hf - E_c = 10 - 5,3 = 4,7\ \mathrm{eV}\ \text{\punts{0,2}} \]
\end{solapartat}
